\documentclass[11pt]{article}
\usepackage{amsmath,amssymb,amsthm}
\usepackage[margin=1in]{geometry}
\newtheorem{theorem}{Theorem}[section]
\newtheorem{lemma}[theorem]{Lemma}
\newtheorem{proposition}[theorem]{Proposition}
\theoremstyle{definition}
\newtheorem{definition}[theorem]{Definition}

\title{Variance bounds for proportions and sums}
\author{}
\date{}

\begin{document}
\maketitle

\begin{abstract}
We record bounds on the variance $p(1-p)/n$ of a sample proportion, a concavity property of its square root, a
uniform bound in $n$, and a bound for the sample variance of a sum of two data vectors.
\end{abstract}

\section{Setting}

Throughout, $n \ge 1$ is an integer and $I$ denotes the unit interval.

\begin{definition}\label{def:vp}
For $p \in \mathbb{R}$, let $V_n(p) = p(1-p)/n$; for $p \in I$ this is the variance of the proportion of successes
in $n$ independent trials with success probability $p$.
\end{definition}

\begin{definition}\label{def:svar}
For $m \ge 2$ and $u \in \mathbb{R}^m$, let $\bar u = \frac1m \sum_{i=1}^m u_i$ and
$s^2(u) = \frac{1}{m-1} \sum_{i=1}^m (u_i - \bar u)^2$.
\end{definition}

\section{Results}

\begin{lemma}\label{lem:quarter}
For every $p \in \mathbb{R}$, $p(1-p) = \frac14 - (p - \frac12)^2$. In particular $V_n(p) \le \frac{1}{4n}$.
\end{lemma}

\begin{proof}
Expanding, $\frac14 - (p - \frac12)^2 = \frac14 - p^2 + p - \frac14 = p(1-p)$. Since $(p - \frac12)^2 \ge 0$ and
$n > 0$, the bound follows.
\end{proof}

\begin{lemma}\label{lem:positive}
For every $p \in I$, $V_n(p) > 0$.
\end{lemma}

\begin{proof}
For $p \in I$ both $p$ and $1 - p$ are positive, and $n > 0$.
\end{proof}

\begin{proposition}\label{prop:concave}
For all $p, q \in I$,
$\sqrt{V_n\bigl(\tfrac{p+q}{2}\bigr)} \ge \frac12 \Bigl(\sqrt{V_n(p)} + \sqrt{V_n(q)}\Bigr)$.
\end{proposition}

\begin{proof}
It is well known that the function $p \mapsto \sqrt{p(1-p)}$ is concave on $I$. The claim is the inequality defining
concavity at the midpoint of $p$ and $q$, divided by $\sqrt{n}$.
\end{proof}

\begin{lemma}\label{lem:linear}
For every $p \in I$, $V_n(p) \le \frac{p}{4n}$.
\end{lemma}

\begin{proof}
By Lemma~\ref{lem:quarter}, $V_n(p) = \frac{p(1-p)}{n} \le \frac{p}{4n}$.
\end{proof}

\begin{proposition}\label{prop:uniform}
There is a constant $C$ such that $V_n(p) \le C/n$ for every $n \ge 1$ and every $p \in \mathbb{R}$.
\end{proposition}

\begin{proof}
By Lemma~\ref{lem:quarter}, $C = \frac14$ works.
\end{proof}

\begin{proposition}\label{prop:sum}
For $m \ge 2$ and $x, y \in \mathbb{R}^m$, $s^2(x + y) \le 2 s^2(x) + 2 s^2(y)$.
\end{proposition}

\begin{proof}
We have $s^2(x + y) = s^2(x) + s^2(y)$, and $s^2(x) + s^2(y) \le 2 s^2(x) + 2 s^2(y)$ because both terms are
nonnegative.
\end{proof}

\end{document}
